% file: graphs.tex

\chapter{Graphs}
\label{chapter:graphs}
\marginnote{Graphs are fundamental data structures used to model pairwise relations between objects. They are widely applicable in various domains such as computer networks, social networks, and transportation systems.}

\section*{Visual Representations}
\begin{figure}[h]
    \centering
    \begin{tikzpicture}
        % Example of different graph types
        \node[draw] at (0,0) {Directed};
        \node[draw] at (2,0) {Undirected};
        \node[draw] at (4,0) {Weighted};
        % Add more visual examples
    \end{tikzpicture}
    \caption{Common Graph Types and Their Representations}
    \label{fig:graph_types}
\end{figure}

\section*{Time and Space Complexity Analysis}
\begin{table}[h]
    \centering
    \begin{tabular}{|l|c|c|c|}
        \hline
        \textbf{Operation} & \textbf{Adjacency Matrix} & \textbf{Adjacency List} & \textbf{Edge List} \\
        \hline
        Add Edge & O(1) & O(1) & O(1) \\
        Remove Edge & O(1) & O(d) & O(E) \\
        Query Edge & O(1) & O(d) & O(E) \\
        Space & O(V²) & O(V + E) & O(E) \\
        \hline
    \end{tabular}
    \caption{Time and Space Complexity Comparison}
    \label{table:complexity}
\end{table}

\section*{Implementation Examples}
\begin{lstlisting}[language=Python]
class Graph:
    def __init__(self, directed=False):
        self.graph = defaultdict(list)
        self.directed = directed
    
    def add_edge(self, u, v, weight=1):
        self.graph[u].append((v, weight))
        if not self.directed:
            self.graph[v].append((u, weight))
    
    def bfs(self, start):
        visited = set()
        queue = deque([start])
        while queue:
            vertex = queue.popleft()
            if vertex not in visited:
                visited.add(vertex)
                queue.extend(v for v, _ in self.graph[vertex] 
                           if v not in visited)
        return visited
\end{lstlisting}

\section*{Common Graph Patterns}
\begin{description}
    \item[Cycle Detection:] 
        \begin{itemize}
            \item Using DFS with color marking
            \item Using Union-Find data structure
            \item Floyd's Cycle-Finding Algorithm
        \end{itemize}
    
    \item[Shortest Path Patterns:] 
        \begin{itemize}
            \item Single-source shortest path
            \item All-pairs shortest path
            \item Bidirectional search
        \end{itemize}
        
    \item[Network Flow:] 
        \begin{itemize}
            \item Ford-Fulkerson Algorithm
            \item Edmonds-Karp Algorithm
            \item Push-Relabel Algorithm
        \end{itemize}
\end{description}

\section*{Advanced Topics}
\begin{description}
    \item[Graph Coloring:] 
        \begin{itemize}
            \item Vertex coloring algorithms
            \item Edge coloring algorithms
            \item Map coloring problems
        \end{itemize}
    
    \item[Maximum Flow:] 
        \begin{itemize}
            \item Maximum flow-minimum cut theorem
            \item Residual graphs
            \item Augmenting paths
        \end{itemize}
        
    \item[Graph Matching:] 
        \begin{itemize}
            \item Bipartite matching
            \item Maximum matching
            \item Perfect matching
        \end{itemize}
\end{description}

\section*{Optimization Techniques}
\begin{description}
    \item[Memory Optimization:] 
        \begin{itemize}
            \item Bit manipulation for graph representation
            \item Compressed sparse row format
            \item Memory-efficient adjacency lists
        \end{itemize}
    
    \item[Performance Optimization:] 
        \begin{itemize}
            \item Parallel graph algorithms
            \item Cache-friendly implementations
            \item Graph partitioning
        \end{itemize}
\end{description}

\section*{Real-World Case Studies}
\begin{description}
    \item[Social Network Analysis:] 
        \begin{itemize}
            \item Friend recommendation systems
            \item Community detection
            \item Influence propagation
        \end{itemize}
    
    \item[Route Planning:] 
        \begin{itemize}
            \item GPS navigation systems
            \item Traffic optimization
            \item Public transportation routing
        \end{itemize}
\end{description}

Graphs consist of a set of nodes (also called vertices) and a set of edges connecting pairs of nodes. They can be categorized based on their properties:

\begin{itemize}
    \item \textbf{Directed vs. Undirected Graphs:} In directed graphs, edges have a direction, indicating a one-way relationship. In undirected graphs, edges have no direction, representing a two-way relationship.
    \index{Directed Graphs}
    \index{Undirected Graphs}
    
    \item \textbf{Weighted vs. Unweighted Graphs:} Weighted graphs have edges associated with weights or costs, while unweighted graphs do not.
    \index{Weighted Graphs}
    \index{Unweighted Graphs}
    
    \item \textbf{Cyclic vs. Acyclic Graphs:} Cyclic graphs contain at least one cycle, whereas acyclic graphs do not.
    \index{Cyclic Graphs}
    \index{Acyclic Graphs}
    
    \item \textbf{Connected vs. Disconnected Graphs:} In connected graphs, there is a path between every pair of nodes. Disconnected graphs consist of multiple connected components.
    \index{Connected Graphs}
    \index{Disconnected Graphs}
\end{itemize}

\subsection*{Graph Representations}
Graphs can be represented in various ways, each with its own advantages:

\begin{itemize}
    \item \textbf{Adjacency Matrix:} A 2D array where each cell \((i, j)\) indicates the presence (and possibly the weight) of an edge between nodes \(i\) and \(j\). This representation allows for quick edge lookups but consumes \(O(n^2)\) space.
    \index{Adjacency Matrix}
    
    \item \textbf{Adjacency List:} An array of lists where each list contains the neighbors of a node. This representation is more space-efficient for sparse graphs and allows for efficient iteration over neighbors.
    \index{Adjacency List}
    
    \item \textbf{Edge List:} A list of all edges in the graph, typically represented as pairs (or triplets if weighted). This representation is simple but less efficient for certain operations.
    \index{Edge List}
\end{itemize}

\subsection*{Graph Traversal Algorithms}
Graph traversal algorithms are essential for exploring the nodes and edges of a graph. The two primary traversal methods are:

\begin{enumerate}
    \item \textbf{Depth-First Search (DFS):} Explores as far as possible along each branch before backtracking. DFS can be implemented using recursion or an explicit stack.
    \index{Depth-First Search (DFS)}
    
    \item \textbf{Breadth-First Search (BFS):} Explores all neighbors at the current depth before moving to nodes at the next depth level. BFS is typically implemented using a queue.
    \index{Breadth-First Search (BFS)}
\end{enumerate}

\subsection*{Advanced Graph Algorithms}
Beyond basic traversal, several advanced algorithms solve complex problems on graphs:

\begin{itemize}
    \item \textbf{Dijkstra's Algorithm:} Finds the shortest path from a single source to all other nodes in a weighted graph with non-negative edge weights.
    \index{Dijkstra's Algorithm}
    
    \item \textbf{Bellman-Ford Algorithm:} Computes shortest paths from a single source in graphs that may have negative edge weights.
    \index{Bellman-Ford Algorithm}
    
    \item \textbf{Floyd-Warshall Algorithm:} Finds shortest paths between all pairs of nodes in a weighted graph.
    \index{Floyd-Warshall Algorithm}
    
    \item \textbf{Kruskal's and Prim's Algorithms:} Used for finding the Minimum Spanning Tree (MST) of a weighted, undirected graph.
    \index{Kruskal's Algorithm}
    \index{Prim's Algorithm}
    
    \item \textbf{Topological Sort:} Orders the nodes of a directed acyclic graph (DAG) such that for every directed edge \( uv \), node \( u \) comes before \( v \) in the ordering.
    \index{Topological Sort}
    
    \item \textbf{Tarjan's Algorithm:} Identifies strongly connected components in a directed graph.
    \index{Tarjan's Algorithm}
\end{itemize}

\subsection*{Graph Theory Concepts}
Understanding fundamental graph theory concepts is crucial for solving graph-related problems:

\begin{itemize}
    \item \textbf{Degree of a Node:} The number of edges incident to a node. In directed graphs, this includes in-degree and out-degree.
    \index{Degree of a Node}
    
    \item \textbf{Path and Cycle:} A path is a sequence of nodes where each consecutive pair is connected by an edge. A cycle is a path that starts and ends at the same node without repeating any edges or nodes.
    \index{Path}
    \index{Cycle}
    
    \item \textbf{Connected Components:} Subsets of nodes where each node is reachable from any other node in the same subset.
    \index{Connected Components}
    
    \item \textbf{Bipartite Graph:} A graph whose nodes can be divided into two disjoint sets such that every edge connects a node from one set to the other.
    \index{Bipartite Graph}
    
    \item \textbf{Graph Isomorphism:} Determines whether two graphs are structurally identical, meaning there's a one-to-one correspondence between their node sets that preserves adjacency.
    \index{Graph Isomorphism}
\end{itemize}

\subsection*{Applications of Graphs}
Graphs are versatile structures used in various real-world applications:

\begin{itemize}
    \item \textbf{Social Networks:} Modeling relationships between individuals.
    \index{Social Networks}
    
    \item \textbf{Computer Networks:} Representing the connectivity of devices.
    \index{Computer Networks}
    
    \item \textbf{Transportation Systems:} Modeling routes and connections between locations.
    \index{Transportation Systems}
    
    \item \textbf{Recommendation Systems:} Suggesting products or content based on user interactions.
    \index{Recommendation Systems}
    
    \item \textbf{Biological Networks:} Representing interactions within biological systems, such as protein-protein interaction networks.
    \index{Biological Networks}
\end{itemize}

\section*{Graph Problems}
The following problems explore various graph concepts and algorithms. Each problem is introduced in its respective section.

\printindex